\chapter{Architettura della compresenza musicale}\label{sec:engineering}

Il centro tecnico del Portale è una relazione fra due persone: suonare
qui, vedere davanti a sé un altro pianista, sentirne la risposta mentre
il proprio gesto sta ancora prendendo forma. Schermo OLED e camera
costruiscono l'ambizione di una presenza quasi speculare; audio
bidirezionale e strumenti locali devono darle continuità. Il collegamento
non è soltanto un accessorio per ricevere lezioni. È una funzione
costitutiva da progettare e verificare, anche quando il desiderio si
spinge fino a un interlocutore distante diecimila chilometri.

La fonte originale integra tastiera MIDI, schermo, calcolo e webcam
\cite[pp.~2--3]{terzi_peeragogy_2016}. La proposta seguente sviluppa
quell'intuizione attraverso due postazioni comunicanti. Non descrive
un sistema già realizzato: componenti, budget e collaudi sono ipotesi
operative. Il funzionamento locale protegge la continuità dello
strumento; la riuscita del progetto connesso si valuta però nella
qualità musicale e visiva dell'incontro fra le due estremità.

\section{Due estremità, un'esperienza condivisa}

Ogni postazione comprende tastiera e pedali, sintesi o acquisizione
audio, mixer, cuffie o diffusori, camera e display. Il pianista ascolta
una combinazione regolabile del proprio suono e del suono remoto.
L'audio viaggia contemporaneamente nei due sensi, \emph{full duplex};
il segnale ricevuto non viene rimandato al mittente. La voce ha un
canale distinto, con controllo esplicito, per poter parlare senza
confondere istruzioni, esecuzione e rumore ambientale.

Il percorso locale tastiera--suono rimane indipendente dal viaggio
in rete. La copia destinata all'altro Portale attraversa acquisizione,
pacchettizzazione, trasporto, buffer di ricezione e riproduzione. Un
secondo percorso porta camera, codifica e immagine al display remoto.
Il piano di controllo coordina invito, accettazione, modalità,
misure e interruzione; il piano di gestione mantiene configurazioni
e software. Nessun aggiornamento può interrompere una sessione attiva.

\begin{figure}[htbp]
\centering
\resizebox{0.98\linewidth}{!}{%
\begin{tikzpicture}[
 box/.style={draw,rounded corners=2pt,text width=3.7cm,minimum height=1cm,align=center,font=\small},
 a/.style={-{Stealth[length=2mm]},thick},
 v/.style={-{Stealth[length=2mm]},dashed}]
\node[box] (aa) {Portale A: gesto / sintesi\\Mixer locale + remoto};
\node[box,right=4cm of aa] (ba) {Portale B: gesto / sintesi\\Mixer locale + remoto};
\node[box,below=1.7cm of aa] (av) {Camera A / display A\\Acquisizione e resa video};
\node[box,below=1.7cm of ba] (bv) {Camera B / display B\\Acquisizione e resa video};
\node[box,below=1.5cm of av] (ac) {Controllo A\\Consenso, stato, diagnostica};
\node[box,below=1.5cm of bv] (bc) {Controllo B\\Consenso, stato, diagnostica};
\draw[a] (aa.north east) to[bend left=15] node[above,font=\small]{Audio A $\rightarrow$ B: $L_{AB}$} (ba.north west);
\draw[a] (ba.south west) to[bend left=15] node[below,font=\small]{Audio B $\rightarrow$ A: $L_{BA}$} (aa.south east);
\draw[v,<->] (av) -- node[above,font=\small]{Video: $V_{AB}$, $V_{BA}$} (bv);
\draw[v,<->] (ac) -- node[above,font=\small]{Stato / invito} (bc);
\draw[v] (ac.west) -- ++(-0.5,0) |- (aa.west);
\draw[v] (bc.east) -- ++(0.5,0) |- (ba.east);
\draw[v] (ac) -- (av);
\draw[v] (bc) -- (bv);
\end{tikzpicture}}
\caption{Architettura proposta a due estremità. Ogni blocco audio
contiene ascolto locale e ricezione remota, senza ritorno del segnale
ricevuto. Audio e video hanno latenze distinte da misurare nei due
sensi; le connessioni possono essere dirette oppure attraversare
un servizio intermedio. Il disegno non implica latenze simmetriche.}
\label{fig:engineering-architecture}
\end{figure}

\section{Trasportare il suono, oppure ricostruirlo da MIDI}

JackTrip offre un riferimento concreto: il progetto documenta
trasmissione audio bidirezionale multicanale non compressa e
implementazioni su più sistemi operativi \cite{net_jacktrip}.
È un candidato da integrare e confrontare, non una soluzione già
collaudata nel Portale. LoLa, sviluppato dal Conservatorio Tartini
con GARR, affronta congiuntamente audio e video per la performance
remota; il manuale 2.0.0 descrive hardware specifico e reti adeguate,
con requisiti elevati nella configurazione video non compressa
\cite{net_lola_manual}. Questi precedenti dimostrano l'esistenza del
campo tecnico, non la fattibilità automatica su una connessione
ordinaria o con una webcam qualsiasi.

L'audio trasmesso conserva il risultato sonoro prodotto dal pianista,
compresi timbro, dinamica ed effetti. Il MIDI trasmette invece eventi:
la postazione remota deve ricostruirne il suono con un sintetizzatore
\cite{eng_midi}. Banche timbriche diverse, curve di velocity, pedali,
polifonia e riverberi possono produrre esperienze differenti. Per una
modalità MIDI condivisa occorre dichiarare strumenti e mappature,
gestire ordinamento e perdita degli eventi, sincronizzare le scelte
ed evitare note sospese. Il minor traffico non elimina il tempo di
propagazione; una previsione di note future altererebbe ciò che
l'altro sta realmente suonando.

Si propone quindi audio come riferimento per la presenza musicale,
con MIDI opzionale per partitura, analisi del gesto e attività
concordate. Inviare entrambi non autorizza a sommare due riproduzioni
della stessa parte: si sceglie esplicitamente quale ascoltare.
Per due canali PCM a 48~kHz e 24 bit, il solo contenuto audio richiede
$48.000\times24\times2=2.304.000$ bit/s per direzione; intestazioni,
ridondanza e video si aggiungono. È un calcolo di dimensionamento,
non una misura di banda sufficiente per una buona esecuzione.

\section{Tre latenze che non vanno confuse}

La latenza \emph{locale} collega gesto e ascolto del proprio strumento.
La latenza \emph{unidirezionale}, $L_{AB}$, collega il gesto in A
all'ascolto in B; $L_{BA}$ descrive il verso opposto. Il \emph{round-trip}
misura un'andata e ritorno definiti: un ping di rete non include
necessariamente acquisizione, sintesi, mixer e trasduttori. Dimezzare
il ping non restituisce una misura del percorso musicale in un senso.

\[
L_{AB}=L_{\mathrm{ingresso},A}+L_{\mathrm{invio},A}
+L_{\mathrm{rete},AB}+L_{\mathrm{buffer},B}
+L_{\mathrm{uscita},B}.
\]

Le voci devono avere confini non sovrapposti; una sintesi software
può sostituire l'acquisizione analogica. Il budget locale proposto
resta mediana entro 15~ms e percentile 95 entro 20~ms, da collaudare.
Questi valori non si trasferiscono al collegamento geografico.
A 48~kHz, 128 campioni occupano 2,67~ms: moltiplicare buffer,
aggiungere codifica e attraversare la rete accumula ulteriore ritardo.

Diecimila chilometri richiedono una formulazione precisa. Usando
come esempio l'indice di gruppo $n_g=1,4682$ riportato da Corning
per una fibra convenzionale a 1550~nm \cite{net_fiber}, la sola
propagazione su 10.000~km di fibra vale:
\[
t_{\mathrm{prop}}=\frac{n_g d}{c}
\simeq\frac{1,4682\times10^{7}}{299.792.458}\ \mathrm{s}
\simeq49\ \mathrm{ms}.
\]
Circa 50~ms in un senso, dunque circa 100~ms per due percorsi
analoghi, prima di apparati, code e conversioni. Il tracciato reale
può superare la distanza geografica; questa è una stima inferiore
per l'ipotesi di fibra dichiarata, non una misura di una rotta.
Un collegamento più capiente non rimuove tale contributo.

Non ne segue che ogni esperienza musicale a quella distanza sia
priva di valore. Segue che la stessa simultaneità fisica non offre
necessariamente lo stesso accoppiamento ritmico di due pianisti
nella medesima stanza. Le prove devono distinguere repertorio,
tempo, ruolo e strategia di ascolto. Non si promette una soglia
universale oltre la quale tutta la musica smette di funzionare.

\section{Jitter, perdita e ascolto senza ritorni}

La variabilità degli arrivi, o jitter, costringe a scegliere fra
attesa e rischio di interruzione. Un buffer maggiore può rendere
continuo il suono aggiungendo ritardo; JackTrip documenta questo
compromesso e distingue diversi contributi nella diagnostica
\cite{net_jacktrip_latency}. Il test deve registrare impostazioni,
ritardi e artefatti: una barra verde non basta a descrivere
l'esperienza. Perdita casuale e raffiche di pacchetti mancanti
vanno provate separatamente.

La configurazione iniziale proposta impiega Ethernet e ascolto
cablato. Ogni estremità ha una propria coda di ricezione, poiché
le condizioni possono essere asimmetriche. Eventuali tecniche di
ridondanza o ricostruzione dei campioni devono essere dichiarate:
mascherare un'interruzione non equivale a recuperare con certezza
l'esecuzione originale. Durante un cambiamento di buffer, entrambi
vedono che la modalità sonora è cambiata. Il sistema non continua
a qualificare come stabile una sessione che sta degradando.

Le cuffie semplificano la prima configurazione full duplex. Con
diffusori e microfono aperto occorre controllare rientri acustici,
eco e feedback. La cancellazione d'eco pensata per il parlato
può alterare suoni sostenuti e sovrapposti: va valutata musicalmente,
non attivata come garanzia generica. Il mixer invia soltanto fonti
locali e mantiene regolazioni indipendenti per sé, interlocutore
e voce. La perdita di un dispositivo silenzia il percorso guasto;
il ripristino non riapre automaticamente microfoni o livelli elevati.

\section{Il quasi specchio: camera, sguardo e tempo audiovisivo}

Lo schermo OLED è una scelta espressiva da verificare insieme a
luminosità, riflessi, dimensione percepita e input lag. La velocità
del pannello non equivale alla latenza camera--display. Esposizione,
acquisizione, codifica, rete, decodifica e visualizzazione possono
ritardare l'immagine. A 60 immagini al secondo un intervallo dura
16,7~ms, ma l'intera catena può richiederne diversi. LoLa è
particolarmente pertinente proprio perché considera acquisizione
e resa audiovisiva come un sistema \cite{net_lola_manual}.

La camera principale dovrebbe avvicinarsi all'asse dello sguardo;
un'inquadratura aggiuntiva delle mani può servire all'osservazione
tecnica. La scelta va prototipata con persone di stature differenti,
sedute e in piedi: guardare il volto sullo schermo non significa
guardare nella camera. L'immagine quasi speculare è un'ipotesi
percettiva, non una giustificazione per invertire automaticamente
mani e tastiera. Orientamento, scala e anteprima locale devono
restare chiari al partecipante.

Si misura separatamente lo scarto $\Delta_{AV}=V_{AB}-L_{AB}$,
riferendo entrambe le catene allo stesso evento fisico e alla stessa
origine temporale. Con riferimenti diversi la sottrazione non descrive
lo scarto percepito fra gesto e suono.
Ritardare l'audio per allinearlo a un video lento può peggiorare
l'interazione musicale. In modalità ensemble si privilegia la
catena audio e si esplicita l'eventuale scarto; in una dimostrazione
si può concordare un diverso compromesso. Nessuna correzione
può mostrare un gesto futuro: la sincronizzazione va ottenuta
riducendo ritardi o introducendo attese note, non inventando
movimenti del pianista.

\section{Pratiche connesse e piattaforma di studio}

Il Portale propone tre regimi espliciti. Nell'\emph{ensemble
connesso} entrambi suonano contemporaneamente, con attività e
condizioni tecniche annotate; la valutazione comprende stabilità
del tempo, ascolto reciproco e libertà espressiva. Nell'\emph{alternanza}
uno propone e l'altro risponde, senza far passare la conversazione
per un duetto simultaneo. Nel lavoro \emph{asincrono} una traccia
viene ripresa successivamente, con provenienza e tempo dichiarati.
Sono pratiche musicali proprie, non livelli di dignità del progetto.

La sperimentazione a grande distanza può includere parti sostenute,
entrate concordate, improvvisazioni che incorporano il ritardo,
un riferimento temporale condiviso o ruoli ritmici asimmetrici.
Queste sono ipotesi compositive da provare, non soluzioni universali.
Un metronomo sincronizzato aiuta a condividere un riferimento ma
non fa arrivare istantaneamente la risposta dell'altro. Il desiderio
di suonare insieme resta la domanda che guida il confronto.

La piattaforma di studio può essere propria oppure collegare
risorse esterne. Deve distinguere materiale curato, sessione dal
vivo e registrazione, consentire scelta di repertorio e scambio
di annotazioni, dichiarare licenze e disponibilità offline.
L'integrazione con un servizio esistente richiede verifica di
accessibilità, esportazione, costi e trattamento dei dati. Non è
necessario costruire un nuovo sistema di videoconferenza per
rendere originale l'esperienza del Portale.

Un aiuto AI può spiegare o proporre un esercizio fuori dal percorso
audio in tempo reale. Non sostituisce l'interlocutore con una
simulazione non dichiarata, non attribuisce note previste al
pianista remoto e non valuta automaticamente l'apprendimento.
Risposte tardive vengono eliminate alla chiusura della sessione.

\section{Invito, consenso e guasto leggibile}

La macchina di stati proposta attraversa \emph{disponibile},
\emph{invito}, \emph{accettazione reciproca}, \emph{prova audio-video},
\emph{sessione}, \emph{degradazione}, \emph{chiusura} e \emph{ripristino}.
Prima dell'accettazione non vengono trasmessi audio o video.
La camera può essere attivata su scelta per la sola anteprima locale,
chiaramente distinta dalla trasmissione; il microfono resta chiuso
finché non viene abilitato per la prova concordata. La prova rende udibili i livelli, verifica la
presenza del partner e mostra la modalità concordata. Entrambe
le persone possono interrompere immediatamente la trasmissione.

Una perdita di rete viene annunciata: nessun filmato in loop
continua a fingere presenza. Il nucleo locale resta suonabile,
ma una ripresa del collegamento richiede conferma, evitando di
riaprire una stanza privata senza preavviso. Passare da ensemble
ad alternanza richiede una scelta condivisa; disattivare la rete
non equivale ad aver completato l'esperienza promessa.

Le sessioni non sono registrate per impostazione iniziale. Un'eventuale
registrazione o trasmissione pubblica richiede un accordo distinto,
con indicazione persistente su entrambi i lati. Inviti temporanei,
autenticazione delle estremità e protezione del trasporto devono
essere verificati nella tecnologia selezionata: un protocollo audio
veloce non offre automaticamente queste proprietà. Le porte di
amministrazione restano separate; i log tecnici escludono contenuti
musicali, immagini e conversazioni salvo ricerca specificamente
concordata.

\section{Collaudo a due estremità e decisioni di progetto}

Il banco iniziale confronta due postazioni sulla stessa rete,
poi una tratta geografica dichiarata e infine condizioni più
impegnative. Si conservano versione software, interfacce, campionamento,
buffer, codec, risoluzione video, instradamento noto e orario della
prova. Un clock condiviso sufficientemente accurato permette
misure unidirezionali; in sua assenza si dichiarano incertezza
ed eventuale sola misura di andata e ritorno. La sincronizzazione
ordinaria dell'orologio non viene presunta precisa al millisecondo.

\begin{table}[htbp]
\centering\small
\begin{tabularx}{\linewidth}{@{}lX@{}}
\toprule
Prova & Evidenza richiesta prima di dichiarare la funzione verificata \\
\midrule
N1 Locale & Gesto--suono in A e B: mediana, p95, massimo, incertezza; confronto col budget locale. \\
N2 Trasporto & $L_{AB}$ e $L_{BA}$, oppure RTT esplicitamente qualificato; non usare ping come misura completa. \\
N3 Continuità & Jitter, perdita e raffiche simulate: artefatti, cambi di buffer, tempi di recupero. \\
N4 Presenza & Camera--display e scarto audio-video nei due sensi; scala, orientamento e sguardo valutati con persone. \\
N5 Musica & Stesso compito in presenza e connesso: andamento temporale, interruzioni e giudizi dei partecipanti. \\
N6 Controllo & Invito rifiutato, revoca, rete caduta, riconnessione: nessuna ripresa non autorizzata. \\
N7 Isolamento & Arresto UI, cavo MIDI scollegato, audio perso: nessuna nota sospesa e guasto riconoscibile. \\
N8 Cura & Ripristino da configurazione versionata e prova prolungata; log privi di contenuti personali non previsti. \\
\bottomrule
\end{tabularx}
\caption{Matrice di collaudo proposta. Nessuna prova è dichiarata
eseguita sul Portale. Le condizioni musicali fanno parte del risultato.}
\end{table}

Le misure tecniche vanno accompagnate da esercizi ripetuti e
resoconti dei due pianisti, evitando di trasformare un'esecuzione
riuscita in garanzia per ogni brano. Il software deve essere
ricostruibile con licenze, dipendenze e configurazioni documentate;
la manutenzione comprende anche camera, cablaggi, ventilazione,
limiti d'ascolto e protezioni fisiche. Un secondo operatore deve
poter ricreare la sessione, non soltanto guardare un video del test.

La sfida non è far sparire la distanza con una dichiarazione.
È costruire una relazione in cui vedere, suonare e rispondere
restino leggibili alle due persone. Il Portale diventa tecnicamente
credibile quando sa mostrare le condizioni nelle quali quella
presenza riesce, e dare forma musicale anche alle distanze che
non può cancellare.
